\textbf{Proof.}
Put \(\varepsilon=\varepsilon_{N_0,K_T,K_D}(\alpha)\), and let \(Z_r\), \(\mathfrak b_r\), and \(Q_r\) be the processes in the corrected linearization lemma. At every sufficiently large triple, that lemma gives
\[
 \sup_P P\{\|Z_r\|_{\infty,\mathcal G}>c_r^\xi+\omega_r/2\}
 \le\alpha-\varepsilon/4,
 \qquad
 \sup_P P\{\|Q_r\|_\infty>\omega_r/2\}\le\varepsilon/4,
\]
and \(\|\mathfrak b_r\|_\infty\le\mathbf 1\{r\ge2\}\eta_D(K_D)\). A union bound therefore gives, uniformly in \(P\),
\[
 P\!\left\{
 \max_{g\in\mathcal G}|\widehat\theta_r(g)-\theta(g)|
 \le c_r^\xi+\mathbf 1\{r\ge2\}\eta_D(K_D)+\omega_r
 \right\}\ge1-\alpha. \tag{1}
\]
This finite-tail statement uses the positive multiplier-quantile slack to pay for the channel-separated coupling error and the centered nonlinear remainder through the explicit buffer \(\omega_r\).

Set \(q=\min\{\beta+1,2\}\). The interpolation proposition and grid condition give
\[
 \|\widehat\theta_r-\theta\|_{\infty,\mathcal I}
 \le\max_{g\in\mathcal G}|\widehat\theta_r(g)-\theta(g)|+c_{\mathrm{int}}h^q,
 \qquad h^q\le\rho_r. \tag{2}
\]
Projection of a real number onto \([0,1]\) cannot increase its distance from \(\theta(g)\in[0,1]\). Linear interpolation of the projected grid values has its maximum absolute deviation from the linear interpolant of \(\theta\) at a grid endpoint, and \(\theta(0)=1\). Thus (1) and (2) imply simultaneous containment of \(\theta\) by the corrected endpoints at every \(t\in\mathcal I\). The probability is at least \(1-\alpha\), which is stronger than \(1-\alpha-\varepsilon\). Hence \(\widehat B_r\in\mathcal H_r(N_0,K_T,K_D;\alpha)\).

Clipping endpoints to \([0,1]\) can only reduce width, so
\[
 w(\widehat B_r)
 \le2c_r^\xi+2c_{\mathrm{int}}h^q
 +2\mathbf 1\{r\ge2\}\eta_D(K_D)+2\omega_r. \tag{3}
\]
For all sufficiently large triples, \(q_{N_0,K_T,K_D}\le1-\alpha/2\). The bounded martingale envelopes and the entropy calculation used in the linearization proof give a uniform conditional maximal second moment of order \(\rho_r^2\); Markov's inequality at the fixed tail probability \(\alpha/2\) therefore yields
\[
 \sup_{P\in\mathfrak P_r}E_Pc_r^\xi\le C\rho_r. \tag{4}
\]
Also \(h^q\le\rho_r\), and the channel-by-channel calculation in the linearization lemma gives \(\omega_r=o(\rho_r)\) uniformly on every three-dimensional tail. Taking expectations in (3), using (4), and increasing \(a_0\) if necessary yield
\[
 \sup_{P\in\mathfrak P_r}E_Pw(\widehat B_r)
 \le C\{\rho_r+\mathbf 1\{r\ge2\}\eta_D(K_D)\}
 =C\bar\rho_r. \tag{5}
\]
At \(r=1\), the risk-exhaustion term is absent by definition. If \(r\ge2\) and \(\eta_D(K_D)=O(\rho_r)\), then \(\bar\rho_r=O(\rho_r)\). No step compares the relative growth of \(K_T\), \(N_0\), and \(K_D\). \(\square\)